Computer vision models have moved rapidly from research prototypes into production pipelines
and edge deployments, powering applications from autonomous inspection robots to real-time
video analytics at the network edge.

\paragraph{Problem.}
General-purpose model-serving systems such as NVIDIA Triton Inference Server~\cite{nvidia2023triton}
and TorchServe~\cite{pytorch2023torchserve} are architected for cloud data centers: they
require Kubernetes or Docker orchestration, expose complex gRPC configuration surfaces, and
carry Python or JVM runtimes that can exceed several gigabytes of disk and memory overhead.
At the other extreme, deploying a model by running a bare PyTorch script provides neither
a stable REST API, nor lifecycle management (loading, unloading, VRAM accounting), nor the
ability to serve multiple models through a single endpoint.
Compounding the infrastructure problem is a license minefield: the most widely used
detection models---the YOLO series---are released under the GNU Affero General Public
License v3.0 (AGPL-3.0)~\cite{ultralytics2023yolo,agpllicense}, which requires any software
that incorporates or exposes YOLO as a network service to publish its entire source code
under AGPL-3.0, foreclosing use in closed-source or commercial products without a
paid commercial license.

\paragraph{Gap.}
Despite the wide variety of available tools, no existing solution simultaneously provides
(a)~a complete CV serving stack---covering detection, segmentation, open-vocabulary
detection, classification, depth estimation, and embedding---in a single deployable unit;
(b)~a guarantee that every included model is permissively licensed (Apache-2.0, MIT, or
BSD) so that downstream products can use it without license encumbrance; and
(c)~a single self-contained binary that runs on edge hardware without requiring a Python
interpreter, GPU driver stack beyond what ONNX Runtime needs, or container runtime.

\paragraph{VisionServe.}
We present VisionServe, a lightweight inference server written in Go that is designed to be
the ``Ollama for Computer Vision''~\cite{ollama2023}: one binary, pulled from a local model
registry, serving any of 16 registered models over a REST API with no Python dependency.
All inference is delegated to ONNX Runtime~\cite{onnxruntime} via a thin cgo binding,
enabling automatic fallback across execution providers (TensorRT $\to$ CUDA $\to$ CoreML
$\to$ DirectML $\to$ OpenVINO $\to$ CPU) without changes to model code.
License safety is enforced structurally: the model registry parser validates each model's
declared license against an Apache-2.0/MIT/BSD allowlist at server startup, refusing to
load any model that does not pass.

\paragraph{Contributions.}
This paper makes the following contributions:

\begin{itemize}
  \item \textbf{A license-enforcing model loader (primary).} To our knowledge the first
        inference runtime to enforce a permissive-license allowlist (Apache-2.0/MIT/BSD)
        \emph{at model-load time}, refusing non-permissive models by construction, and binding
        the declared license to specific weight bytes (sha256) and an audited
        \texttt{source\_url}, entirely offline in a single edge binary.

  \item \textbf{A novelty boundary and threat model.} We position the gate against build-time
        SCA scanners, deploy-time admission controllers, and model signing---none of which
        enforces a license policy at model-load time---and give an honest threat model
        exercised by a reproducible adversarial demo (an AGPL-declared model, a weight-hash
        mismatch, and an unaudited source are each refused).

  \item \textbf{A CV-specific serving-overhead crossover characterization (measured).} On
        identical ONNX, a single-binary Go server scales to $\sim$$8\times$ a GIL-bound
        FastAPI+ONNX~Runtime baseline's peak goodput (695 vs.\ 85~rps for MobileNetV3-Small at
        concurrency~64), and we report honestly where it loses ($\sim$78\,ms vs.\
        $\sim$15\,ms per request at concurrency~1).

  \item \textbf{A latency decomposition localizing the overhead} to pure-Go JPEG decode
        ($\sim$19\,ms) and resize/normalize ($\sim$3.6\,ms), not the GPU (ORT~Run $<$1\,ms on
        an A6000)---the CV-specific instance of non-inference overhead dominating small-DNN
        serving~\cite{abouelhamayed2024beyond}.

  \item \textbf{A reproducible end-to-end evaluation:} served accuracy (export and serving
        preserve ImageNet Top-1 to within 0.5\%), throughput--latency behavior across
        concurrency, and multi-model footprint (one Go process serving four models uses
        917\,MB RSS / 920\,MB VRAM vs.\ 1632\,MB / 1384\,MB for four FastAPI+ORT processes),
        plus a single unified result schema and Python/JavaScript client libraries across all
        16 registered models (12 benchmarked) on an NVIDIA RTX~A6000.
\end{itemize}

\paragraph{Paper outline.}
The remainder of this paper is organized as follows: Section~\ref{sec:related} surveys
related work in model serving and CV tooling; Section~\ref{sec:design} describes the
system architecture and key design decisions; Section~\ref{sec:catalog} presents the model
catalog; Section~\ref{sec:eval} reports benchmark results; Section~\ref{sec:applications}
illustrates representative use cases; and Sections~\ref{sec:limitations}
and~\ref{sec:conclusion} discuss limitations, future directions, and conclusions.
